\documentclass[12pt]{article}
\usepackage[utf8]{inputenc}
\usepackage[T2A]{fontenc}
\usepackage[russian]{babel}
\usepackage{amsmath}
\usepackage{amssymb}
\usepackage{graphicx}
\usepackage{geometry}
\geometry{margin=2cm}

\title{Уравнение Камерлинга-Оннеса для реального газа}

\date{}

\begin{document}

\maketitle

\section*{Введение}
Уравнение Камерлинга-Оннеса представляет собой один из наиболее универсальных подходов к описанию реальных газов. В отличие от полуэмпирических уравнений состояния (например, Ван-дер-Ваальса), оно основывается на строгом статистическом подходе и учитывает межмолекулярные взаимодействия через последовательное разложение давления в ряд по степеням плотности газа.

\section{Основные положения вириального разложения}
Для идеального газа уравнение состояния имеет вид:
$$
PV = RT
$$
где $P$ — давление, $V$ — молярный объем, $T$ — температура, $R$ — универсальная газовая постоянная.

Для реального газа необходимо учитывать поправки, связанные с:
\begin{enumerate}
    \item Конечным размером молекул
    \item Межмолекулярными силами притяжения и отталкивания
\end{enumerate}

Эти эффекты можно систематически учесть, представив давление в виде вириального разложения:
\[
PV = RT \left(1 + \frac{B_2(T)}{V} + \frac{B_3(T)}{V^2} + \frac{B_4(T)}{V^3} + \dots \right)
\]
где $B_n(T)$ — вириальные коэффициенты $n$-го порядка.

\section{Вывод уравнения}
Вириальное разложение основывается на статистической механике и теории возмущений. Рассмотрим ключевые этапы вывода:

\subsection{1. Канонический ансамбль}
Свободная энергия Гельмгольца $F$ выражается через статистический интеграл:
$$
F = -kT \ln Z
$$
где $Z$ — статистическая сумма системы, $k$ — постоянная Больцмана.

\subsection{2. Разложение по плотности}
Разложим давление в ряд по степеням плотности $\rho = 1/V$:
\[
\frac{PV}{RT} = 1 + B_2(T)\rho + B_3(T)\rho^2 + B_4(T)\rho^3 + \dots
\]

\subsection{3. Связь с межмолекулярными силами}
Вириальные коэффициенты выражаются через интегралы от функций распределения:
\[
B_2(T) = -2\pi \int_0^\infty \left( e^{-U(r)/kT} - 1 \right) r^2 dr
\]
где $U(r)$ — потенциал межмолекулярного взаимодействия.

\section{Физический смысл вириальных коэффициентов}

    
- $B_2(T)$ — учитывает парные взаимодействия молекул
    
- $B_3(T)$ — тройные взаимодействия
    
- $B_4(T)$ — четверные взаимодействия и т.д.


При низких плотностях ($V \to \infty$) достаточно учитывать только $B_2(T)$. При увеличении плотности необходимы более высокие члены разложения.

\section{Примеры вириальных коэффициентов}
Для конкретных потенциалов межмолекулярного взаимодействия можно вычислить вириальные коэффициенты аналитически:

    
- Потенциал Леннарда-Джонса:
    $$
    U(r) = 4\epsilon \left[ \left( \frac{\sigma}{r} \right)^{12} - \left( \frac{\sigma}{r} \right)^6 \right]
    $$
    
- Для потенциала твердых сфер ($U(r)=\infty$ при $r<d$, $U(r)=0$ при $r\geq d$):
    $$
    B_2(T) = \frac{2}{3}\pi d^3
    $$


\section{Сравнение с другими уравнениями состояния}
\begin{center}
\begin{tabular}{|c|c|}
\hline
Модель & Уравнение \\
\hline
Идеальный газ & $PV = RT$ \\
\hline
Ван-дер-Ваальс & $\left(P + \frac{a}{V^2}\right)(V - b) = RT$ \\
\hline
Камерлинг-Оннес & $PV = RT\left(1 + \frac{B_2}{V} + \frac{B_3}{V^2} + \dots\right)$ \\
\hline
\end{tabular}
\end{center}

\section{Применение уравнения}
Уравнение Камерлинга-Оннеса используется:

    
- Для точных термодинамических расчетов при умеренных давлениях
    
- При построении стандартных шкал температуры
    
- Для анализа критических явлений
    
- В исследованиях сверхтекучих систем


\section*{Заключение}
Вириальное разложение Камерлинга-Оннеса предоставляет универсальный метод описания реальных газов, основываясь на фундаментальных принципах статистической физики. Его преимущество заключается в возможности последовательного учета межмолекулярных взаимодействий с заданной точностью. Однако при высоких плотностях сходимость ряда замедляется, что требует применения альтернативных подходов.

\end{document}